\documentclass[11pt]{article}
\usepackage[margin=1in]{geometry}
\usepackage[hidelinks]{hyperref}
\setlength{\parindent}{0pt}
\setlength{\parskip}{6pt}
\begin{document}
\begin{center}
{\Large\bfseries Cover Letter\par}
\smallskip
\emph{Auditing Hidden Reasoning Charges in Commercial LLM APIs through Token-Usage Prediction}
\end{center}

Dear Dr. Liu,

Thank you for the opportunity to revise our manuscript (ID b52fcf7f-1f82-46db-8e40-8e72a1746bdf). The manuscript was submitted to the \emph{Securing the Core: Preventing Unauthorized Use and Manipulation of Large Language Models} Collection in \emph{npj Artificial Intelligence}. We have revised the manuscript and provided point-by-point responses to the reviewers in the accompanying rebuttal letter. Our response to your editorial comment is below.

\section*{Editorial Comment}
\begin{quote}\small
All reviewers found the research problem interesting and acknowledged the strengths of the manuscript. However, they also raised several concerns, most of which appear to be addressable through revision. Therefore, I invite the authors to address the reviewers' comments and provide a response before the manuscript can be considered further.
\par
Please ensure the results are accurately reported, any overstated conclusions are rewritten and the limitations of the work fully explained.
\end{quote}

\section*{Response to the Editor}

\paragraph{1. Scope and strength of the conclusions.}
PALACE turns user-observable prompts and final answers into independent usage estimates that can be aggregated to monitor reported reasoning-token charges. The revision experiments quantify this capability on held-out workloads. We distinguish count prediction (the tolerance sweep in Figure~3) from detection of inflated reports (Supplementary Section~E). The Abstract, the Results subsection on simulated reporting discrepancies, and the Discussion report the workload-level detection result.

\paragraph{2. Detection rate for simulated inflation.}
We evaluated simulated overreporting directly on held-out data. For the PALACE domain-specific 1.5B auditor, 30\% overreporting was detected in 92.9\% of 50-request batches at 4.9\% observed false-positive rate (FPR), and in 99.7\% of 100-request batches at 5.2\% FPR. The batches mixed the four domains uniformly. At 10\% overreporting, the first tested size exceeding 90\% detection was 500 requests (96.5\% detection, 8.4\% FPR). Supplementary Section~E reports the workload-size sweep and observed PALACE FPR for each tested batch size.

\paragraph{3. Limitations.}
The detector uses reference counts matched to the audited model and configuration and sets its threshold on clean calibration traffic. Supplementary Section~E quantifies how reasoning-length and domain mix shifts affect the false-positive rate. In a workload with 70\% coding requests, domain-wise calibration reduced the 100-request false-positive rate from 25.9\% to 3.5\%. The evaluation uses trace-derived references and simulated reporting changes.

Sincerely,\\
The authors
\end{document}
